% Generated draft: requires mathematical and source review.
\documentclass[11pt]{article}
\usepackage[margin=25mm]{geometry}
\usepackage{fontspec}
\setmainfont{Latin Modern Roman}
\usepackage{amsmath,amssymb,mathtools,unicode-math}
\setmathfont{Latin Modern Math}
\usepackage{enumitem,needspace,xurl,hyperref}
\setlength{\parindent}{0pt}
\setlength{\parskip}{5pt}
\setlength{\emergencystretch}{3em}
\newcommand{\N}{\mathbb{N}}
\newcommand{\Z}{\mathbb{Z}}
\newcommand{\Q}{\mathbb{Q}}
\newcommand{\R}{\mathbb{R}}
\newcommand{\C}{\mathbb{C}}
\newcommand{\F}{\mathbb{F}}
\newcommand{\A}{\mathbb{A}}
\newcommand{\PP}{\mathbb{P}}
\newcommand{\E}{\mathbb{E}}
\newcommand{\eps}{\varepsilon}
\newcommand{\dd}{\,\mathrm{d}}
\begin{document}
\section*{UnsolvedMath ID: 700824}
\textbf{Source problem code:} catalogue entry
\subsection{Definitions and mathematical statement}
Let \(G\) be an arbitrary group, with no finiteness or finite-generation assumption. Write \([x,y]=x^{-1}y^{-1}xy\) and let
\[
G'=\langle [x,y]\mid x,y\in G\rangle
\]
be its derived subgroup. Assume that
\[
\forall x,y\in G,\qquad [x,y]^3=1.
\]
A group \(H\) has finite exponent if there exists an integer \(n\geq 1\) such that \(h^n=1\) for every \(h\in H\). Define the derived series by \(G^{(0)}=G\) and \(G^{(i+1)}=(G^{(i)})'\) for integers \(i\geq 0\). The group \(G\) is soluble if \(G^{(d)}=\{1\}\) for some integer \(d\geq 0\).

\textbf{Questions (R. Brandl, 1992).} For every group \(G\) satisfying the displayed commutator identity, must each of the following hold?
\begin{enumerate}
\item There exists an integer \(n\geq 1\) such that \(g^n=1\) for every \(g\in G'\).
\item There exists an integer \(d\geq 0\) such that \(G^{(d)}=\{1\}\).
\end{enumerate}
The bounds \(n\) and \(d\) may depend on \(G\); the supplied questions do not explicitly request uniform bounds over all such groups.

The supplied statement attributes to N. D. Gupta and N. S. Mendelssohn (1967) the result that \(G'\) is a \(3\)-group. Here this means that every element of \(G'\) has finite order equal to a power of \(3\), without asserting that \(G'\) is finite or that these orders have a common bound. This attributed result is supplied background, not independently verified here; its full bibliographic reference is absent from the record.
\subsection{Short English statement}
R. Brandl asks whether every group in which every commutator has cube equal to the identity has a derived subgroup of finite exponent, and whether every such group is soluble. The supplied background attributes to Gupta and Mendelssohn (1967) the fact that every element of the derived subgroup has order a power of \(3\).
\subsection{Sources}
\begin{itemize}
\item Local UnsolvedMath registry, supplied complete record: permanent ID \(700824\), exact source code \texttt{catalogue entry}, title \emph{Kourovka Notebook Problem 12.4}. The registry labels the problem \texttt{open}; that label is supplied curation, not an independently verified status assessment.
\item \emph{The Kourovka Notebook (Unsolved Problems in Group Theory)}, original problem 12.4, 12th issue (1992); proposed by R. Brandl. The supplied background identifies the later source as the 21st issue edition (2026), \texttt{arXiv:1401.0300v48}. The supplied collection description names E. I. Khukhro and V. D. Mazurov as editors.
\item N. D. Gupta and N. S. Mendelssohn (1967): the supplied statement attributes to them the result that \(G'\) is a \(3\)-group under the stated hypothesis. No publication title, journal, pages, or direct reference for this result is supplied.
\end{itemize}
\end{document}
